\section{From voltage limits to the joint spectral limit}\label{sec:spectrum}

The four-point voltage theorem identifies the macroscopic part of the
inverse operator.  To deduce spectral convergence, we must also control
functions whose energy is bounded but whose mass could concentrate in a
small region.  The block estimate of Theorem~\ref{thm:block-approximation}
provides this control.  We first turn it into strong compactness in one
Hilbert space, then combine that compactness with the voltage limit to obtain
convergence of the inverse operators in Hilbert--Schmidt norm.

Throughout the first two subsections, work on an almost-sure outcome of a
common representation extraction on which Proposition~\ref{prop:geometry},
Theorem~\ref{thm:block-approximation}, and Theorem~\ref{thm:four-point}
hold simultaneously.  Relabel its indices by \(n\), write \(\mu=\mu_h\),
and retain the two rerootings used for the finite interval cover in
Theorem~\ref{thm:block-approximation}.  All constants and finite starting
indices in this argument may depend on this outcome.  Their asserted
uniformity in later indices and levels is pathwise.

\subsection{Corner time and strong energy compactness}

Let \(H=L^2([0,1],\dd t)\).  For \(r=1,2\), define
\(U_n^{(r)}:L^2(\mu_n)\to H\) by
\[
 (U_n^{(r)}f)(t)=f\bigl(v_n^{(r)}(i)\bigr),
 \qquad \frac{i}{2n}\le t<\frac{i+1}{2n},\quad 0\le i<2n.
\]
Values at \(t=1\) are immaterial.  The exact corner identity
\eqref{eq:corner-pushforward} gives
\[
 \norm{U_n^{(r)}f}_H^2
 =\frac1{2n}\sum_{i=0}^{2n-1}\abs{f(v_n^{(r)}(i))}^2
 =\norm{f}_{L^2(\mu_n)}^2.
\]
Set \(\mathcal H_n=U_n^{(1)}L^2(\mu_n)\).  It is a finite-dimensional subspace of
\(H\) containing the constants.  We conjugate \(B_n\) by \(U_n^{(1)}\)
when it acts on \(\mathcal H_n\), retaining its name, and write
\[
 \mathcal E_n^0(F)=\mathcal E_n^0(f)
 =\ip{F}{B_nF}_H,\qquad F=U_n^{(1)}f\in \mathcal H_n.
\]

The two explorations list the same \(2n\) physical corner occurrences.
Let \(\sigma_n\) be the permutation which matches each occurrence in the
first list to that occurrence in the second.  It preserves its incident
vertex.  For \(t=(i+u)/(2n)\), where \(0\le i<2n\) and \(0\le u<1\), set
\[
 \tau_n(t)=\frac{\sigma_n(i)+u}{2n}.
\]
This is a Lebesgue-measure-preserving bijection up to endpoints, and
\begin{equation}\label{eq:tour-transfer}
 U_n^{(1)}f=(U_n^{(2)}f)\circ\tau_n\quad\text{almost everywhere}.
\end{equation}

\begin{lemma}[Matching the two corner orders]\label{lem:tour-order}
There is a measurable, measure-preserving map \(\tau:[0,1]\to[0,1]\),
defined up to a null set, such that
\[
 \tau_n(t)\longrightarrow\tau(t),\qquad
 \eta_2(\tau(t))=\eta_1(t)
 \quad\text{for almost every }t.
\]
If intervals \(J_n\subset[0,1]\) have endpoints converging to the endpoints
of an interval \(J\), then
\[
 \ind_{J_n}\longrightarrow\ind_J,\qquad
 \ind_{J_n}\circ\tau_n\longrightarrow\ind_J\circ\tau
 \quad\text{in }H.
\]
Finite products and differences of such indicator functions converge in
\(H\) as well.  For either \(r\), the map
\[
 U_r:L^2(\mu)\longrightarrow H,\qquad U_rg=g\circ\eta_r,
\]
is unitary.
\end{lemma}

\begin{proof}
By Proposition~\ref{prop:geometry}, for \(\mu\)-almost every \(z\) the
fiber \(\eta_2^{-1}(\{z\})\) has exactly one element among all times in
\([0,1]\).  The pushforward identity for \(\eta_1\) therefore implies that,
for almost every \(t\), \(\eta_1(t)\) has such a unique second time; call it
\(\tau(t)\).  Fix one of these \(t\), also avoiding the countable union of
step endpoints.  Every subsequence of \(\tau_n(t)\) has a further
convergent subsequence, with some limit \(s\).  The two matched corners have
the same projected vertex.  The uniform tracking in
\eqref{eq:uniform-tour-tracking} and the continuity of \(\eta_2\) give
\[
 \eta_2(s)=\lim_n z_n\bigl(v_n^{(2)}(\sigma_n(\lfloor2nt\rfloor))\bigr)
 =\lim_n z_n\bigl(v_n^{(1)}(\lfloor2nt\rfloor)\bigr)=\eta_1(t)
\]
along this further subsequence.  Singleton fibers force \(s=\tau(t)\).
Thus the entire sequence converges.  The all-times formulation of the
singleton property is what also determines limits that could lie in a null
set of second-tour times.

Extend \(\tau\) arbitrarily on the exceptional null set.  It is measurable
as an almost-everywhere pointwise limit of the measurable \(\tau_n\).  For
every continuous \(g\) on \([0,1]\), bounded convergence yields
\[
 \int_0^1g(\tau(t))\dd t
 =\lim_n\int_0^1g(\tau_n(t))\dd t=\int_0^1g(s)\dd s.
\]
Hence \(\tau\) preserves Lebesgue measure.  Away from the two limiting
endpoints, membership in \(J_n\) eventually agrees with membership in
\(J\).  For the composed indicators the only additional exceptional times
are those for which \(\tau(t)\) is an endpoint; these have measure zero by
the pushforward just proved.  The indicators converge almost everywhere
and are bounded by one, so they converge in \(L^2\).  The assertion for
finite products and differences follows in the same way.

The area pushforward makes \(U_r\) an isometry.  To see surjectivity, let
\(E_r\) be the set of points whose fiber under \(\eta_r\) is a singleton.
For every \(k\), the set of points having two preimages at distance at least
\(1/k\) is the continuous image of a compact subset of \([0,1]^2\), hence
is closed.  Subtracting the union of these sets from the compact image of
\(\eta_r\) shows that \(E_r\) is Borel.  On \(E_r\) the inverse of
\(\eta_r\) is continuous in the relative topology.  Indeed, convergence
of image points and compactness of the time interval show that any limit
of their inverse times belongs to the limiting fiber, which is a singleton.
Choose a Borel representative of any element of \(H\) and compose it with
this inverse on \(E_r\), extending by zero outside \(E_r\).  This constructs
a measurable function of the image whose pullback agrees with the given
time function on \(\eta_r^{-1}(E_r)\).  Since \(\mu(E_r)=1\), the isometry
is onto.
\end{proof}

We now use the finite cover in Theorem~\ref{thm:block-approximation}.
It consists, on the present outcome, of finitely many protected intervals
\(I_{n,\ell}\) in tour order \(r_\ell\in\{1,2\}\), indexed by
\(1\le\ell\le L\).  Each interval is a union of step cells and its duration
converges to a number \(d_\ell>0\).  Every corner is in at least one of
these intervals for all
sufficiently large \(n\).  At level \(j\), write
\(J_{n,\ell,j,b}\), \(0\le b<2^j\), for its partition by nested integer
halves, retaining empty or singleton blocks when necessary.  For each fixed
\(j\), all these blocks are nonempty for large \(n\), and
\begin{equation}\label{eq:block-positive-length}
 \abs{J_{n,\ell,j,b}}\longrightarrow 2^{-j}d_\ell>0.
\end{equation}
This follows by dividing the limiting duration into \(2^j\) nested halves;
the integer-rounding error is \(O(1/n)\) at fixed \(j\).  The absolute
endpoints lie in \([0,1]\), so every subsequence has a further one on which
the endpoints of all \(L\) intervals converge.  Along such a further
subsequence, their fixed-level block endpoints converge as well.

Let \(P_{n,\ell,j}f\) be the block-average function on \(I_{n,\ell}\)
formed from \(U_n^{(r_\ell)}f\).  The theorem supplies one finite level
\(J\), one eventual index \(n_0\), and finite positive constants
\(C_\ell\), such that, simultaneously for \(n\ge n_0\), \(j\ge J\), and all
vertex functions \(f\),
\begin{equation}\label{eq:covered-block-estimate}
 \int_{I_{n,\ell}}
    \abs{U_n^{(r_\ell)}f-P_{n,\ell,j}f}^{2}\dd t
 \le C_\ell 2^{-pj}\mathcal E_n^0(f;F_{n,\ell}),
 \qquad p=\frac7{10}.
\end{equation}
Here \(F_{n,\ell}\) is the edge set of the retained larger graph, a subset
of the full map edge set.  The same \(J,n_0,C_\ell\) work for every finer
level, including the levels at which the partition consists of singletons.
We have taken the maximum of the finitely many thresholds furnished by the
cover theorem.

\begin{lemma}[Finite-rank approximation and energy compactness]
\label{lem:energy-compactness}
On the present extraction there are constants \(C_0,C_1<\infty\) and
linear maps \(Q_{n,j}:\mathcal H_n\to H\), for \(n\ge n_0\) and \(j\ge J\), with
\begin{equation}\label{eq:finite-rank-approximation}
 \rank Q_{n,j}\le C_0 2^j,\qquad
 \norm{F-Q_{n,j}F}_H^2
 \le C_1 2^{-pj}\mathcal E_n^0(F),\quad F\in \mathcal H_n.
\end{equation}
For each fixed \(j\ge J\), whenever \(n_q\to\infty\) and
\(F_q\in\mathcal H_{n_q}\) have bounded norms, the sequence
\(Q_{n_q,j}F_q\) has a strongly convergent subsequence in \(H\).
Consequently, whenever \(n_q\to\infty\) and \(F_q\in\mathcal H_{n_q}\) satisfy
\[
 \sup_q\bigl(\norm{F_q}_H^2+\mathcal E_{n_q}^0(F_q)\bigr)<\infty,
\]
the sequence \(F_q\) has a strongly convergent subsequence in \(H\).
There also exist a finite
\(i_0\) and a constant \(c_*>0\) such that, for all \(n\ge n_0\) and
\(i\ge i_0\),
\begin{equation}\label{eq:raw-eigenvalue-tail}
 \nu_{n,i}\ge c_* i^p.
\end{equation}
The convention \(\nu_{n,i}=+\infty\) for \(i\ge\dim \mathcal H_n\) is used here.
\end{lemma}

\begin{proof}
Put \(\tau_n^{(1)}=\Id\) and \(\tau_n^{(2)}=\tau_n\).  Assign a time to
the first interval that covers it by setting
\[
 \chi_{n,\ell}(t)=
 \ind_{I_{n,\ell}}(\tau_n^{(r_\ell)}(t))
 \prod_{q<\ell}\bigl(1-
       \ind_{I_{n,q}}(\tau_n^{(r_q)}(t))\bigr).
\]
For large \(n\), these functions take values in \(\{0,1\}\) and sum to
one almost everywhere.  For a nonempty original block, let
\[
 \bar f_{n,\ell,j,b}=
 \frac{1}{\abs{J_{n,\ell,j,b}}}
 \int_{J_{n,\ell,j,b}}U_n^{(r_\ell)}f(t)\dd t.
\]
Define the linear map on \(F=U_n^{(1)}f\) by
\begin{equation}\label{eq:priority-approximant}
 Q_{n,j}F=
 \sum_{\ell=1}^L\ \sum_{b:J_{n,\ell,j,b}\ne\varnothing}
 \bar f_{n,\ell,j,b}\,
 \chi_{n,\ell}\,
 \bigl(\ind_{J_{n,\ell,j,b}}\circ\tau_n^{(r_\ell)}\bigr).
\end{equation}
Thus the coefficient is always the average on the full original block;
the priority rule only restricts the indicator on which this coefficient
is used.  There are at most \(L2^j\) summands with fixed indicator
functions, proving the rank bound with \(C_0=L\).

Equation~\eqref{eq:tour-transfer} and the disjoint priority assignment give
\begin{align*}
 \norm{F-Q_{n,j}F}_H^2
 &\le \sum_{\ell=1}^L
       \int_{I_{n,\ell}}
          \abs{U_n^{(r_\ell)}f-P_{n,\ell,j}f}^{2}\dd t \\
 &\le \sum_{\ell=1}^L C_\ell2^{-pj}
                  \mathcal E_n^0(f;F_{n,\ell})
 \le C_1 2^{-pj}\mathcal E_n^0(F),
\end{align*}
where \(C_1=\sum_\ell C_\ell\).  The last inequality uses that every
\(F_{n,\ell}\) is a subset of the original edge set, including its
multiplicities.  The map \(Q_{n,j}\) is allowed to take values outside
\(\mathcal H_n\); only its rank and its approximation error are used.

Fix \(j\) and any sequence of indices and norm-bounded functions as in
the statement.  Relabel its indices by \(n\), and first extract a
subsequence on which the finitely many absolute interval endpoints
converge.  By Cauchy--Schwarz and the isometry of \(U_n^{(r)}\),
\begin{equation}\label{eq:block-coefficient-bound}
 \abs{\bar f_{n,\ell,j,b}}
 \le \abs{J_{n,\ell,j,b}}^{-1/2}\norm{F}_H.
\end{equation}
Equation~\eqref{eq:block-positive-length} bounds these coefficients
uniformly for a norm-bounded sequence.  Lemma~\ref{lem:tour-order} shows
that each of the finitely many indicator factors in
\eqref{eq:priority-approximant}, and hence their product, converges strongly
in \(H\).  After a subsequence the finitely many scalar coefficients also
converge.  This proves the fixed-level compactness.  A priority piece can
have limiting length zero; the bound still holds because the coefficient
was taken on the original block in \eqref{eq:block-coefficient-bound}.

For the asserted energy compactness, start with any sequence of indices
and functions with both norms and energies bounded, again relabelling its
indices by \(n\).  For each successive integer \(j\ge J\), extract a
subsequence on which \(Q_{n,j}F_n\) converges, and use a diagonal
subsequence.  Given \(\varepsilon>0\), choose \(j\) so large that the square
root of the right side of \eqref{eq:finite-rank-approximation} is less than
\(\varepsilon/3\) uniformly along this sequence.  For two sufficiently
late terms, the difference of their level-\(j\) approximants is also less
than \(\varepsilon/3\).  The triangle inequality makes the diagonal
sequence Cauchy in \(H\), proving strong convergence.

Finally let \(i<\dim \mathcal H_n\) be sufficiently large and set
\[
 j(i)=\left\lfloor\log_2\frac{i}{2C_0}\right\rfloor.
\]
Choose \(i_0\) so that \(j(i)\ge J\) for \(i\ge i_0\).  The span of the
first \(i+1\) orthonormal eigenvectors of \(B_n\), including its constant
mode, has dimension greater than \(\rank Q_{n,j(i)}\le i/2\).  It contains
a unit vector \(F\) killed by \(Q_{n,j(i)}\).  Its energy is at most
\(\nu_{n,i}\), so \eqref{eq:finite-rank-approximation} gives
\[
 1\le C_1 2^{-pj(i)}\nu_{n,i},\qquad
 \nu_{n,i}\ge C_1^{-1}\left(\frac{i}{4C_0}\right)^p.
\]
This is \eqref{eq:raw-eigenvalue-tail} with
\(c_*=C_1^{-1}(4C_0)^{-p}\).  Padded eigenvalues are infinite and satisfy
the same bound.
\end{proof}

\subsection{Convergence of the centered inverse kernels}

A voltage difference observes a particular combination of four kernel
values.  For \(k\in L^2([0,1]^2)\), define
\[
 (\mathfrak Dk)(t,t',s,s')
 =k(t,s)-k(t,s')-k(t',s)+k(t',s').
\]
This is a bounded map from \(L^2([0,1]^2)\) to \(L^2([0,1]^4)\), with norm
at most four.  Call \(k\) doubly centered if its integral in either variable
vanishes for almost every value of the other.  These kernels form a closed
subspace, and Fubini's theorem gives
\begin{equation}\label{eq:double-difference-injective}
 \int_0^1\!\int_0^1(\mathfrak Dk)(t,t',s,s')\dd t'\dd s'=k(t,s),
 \qquad \norm{\mathfrak Dk}_{L^2([0,1]^4)}=2\norm{k}_{L^2([0,1]^2)}
\end{equation}
for every doubly centered \(k\).  For the norm identity, expand the square:
the four squared terms each have norm \(\norm{k}_2^2\), and every cross
term vanishes by one of the two centering identities.  In particular, the
four-point difference determines a doubly centered \(L^2\) kernel uniquely.

The general form-convergence theory of
\cite[Definitions~2.11--2.13 and Theorem~2.4]{KuwaeShioya2003} likewise
separates asymptotic compactness of bounded-energy families from
identification of the limiting form.  The next lemma proves the implication
needed here directly from the four-point observation.  It also explains why
a bound for high eigenvalue indices alone does not control the largest
inverse eigenvalue.

\begin{lemma}[A compactness criterion for inverse kernels]
\label{lem:inverse-kernel-compactness}
Let \(H=L^2([0,1])\), and let \(\mathcal H_n\subset H\) be finite-dimensional
subspaces containing the constants.  Let \(B_n\) be nonnegative
self-adjoint operators on \(\mathcal H_n\), with kernel exactly the constants.
Write \(\mathfrak e_n(F)=\ip{F}{B_nF}_H\), list their eigenvalues as
\(0=\nu_{n,0}<\nu_{n,1}\le\cdots\), and pad this list by \(+\infty\) past
\(\dim \mathcal H_n-1\).  Let \(K_n\) be the inverse of \(B_n\) on
\(\mathcal H_n\cap\ind^\perp\), extended by zero on constants and on \(\mathcal H_n^\perp\),
and let \(k_n\in L^2([0,1]^2)\) be its kernel.  Suppose that:
\begin{enumerate}[label=\textup{(\roman*)}]
\item whenever \(n_q\to\infty\) and \(F_q\in \mathcal H_{n_q}\) satisfy
\(\sup_q(\norm{F_q}_H^2+\mathfrak e_{n_q}(F_q))<\infty\), some subsequence
of \(F_q\) converges strongly in \(H\);
\item for some \(p>1/2\), \(c>0\), and finite \(i_0,n_0\), one has
\(\nu_{n,i}\ge c i^p\) for every \(n\ge n_0\) and \(i\ge i_0\);
\item there is a doubly centered \(k\in L^2([0,1]^2)\) for which
\(\mathfrak Dk_n\to\mathfrak Dk\) almost everywhere on \([0,1]^4\).
\end{enumerate}
Then the operator norms \(\norm{K_n}_{\mathrm{op}}\) are bounded for large
\(n\), and \(k_n\to k\) in \(L^2([0,1]^2)\).
\end{lemma}

\begin{proof}
The nonzero eigenvalues of \(K_n\), in decreasing order and with zeros
appended, are \(\kappa_{n,i}=\nu_{n,i}^{-1}\), \(i\ge1\), where
\((+\infty)^{-1}=0\).  Its kernel is doubly centered because \(K_n\) is
self-adjoint and kills constants.  For \(N\ge i_0\), assumption (ii) gives
the uniform Hilbert--Schmidt tail bound
\begin{equation}\label{eq:inverse-square-tail}
 \sum_{i>N}\kappa_{n,i}^2
 \le c^{-2}\sum_{i>N}i^{-2p}=:r_N^2,
 \qquad n\ge n_0,
\end{equation}
where \(r_N\to0\) because \(2p>1\).

Suppose first that \(M_n=\norm{K_n}_{\mathrm{op}}\) is unbounded.  On a
subsequence it tends to infinity, so we may assume \(M_n\ge1\).  Choose
orthonormal eigenvectors \(e_{n,i}\) for the positive eigenvalues and put
\(\alpha_{n,i}=\kappa_{n,i}/M_n\).  Every \(\alpha_{n,i}\) belongs to
\([0,1]\).  By a diagonal extraction, assume these numbers converge to
\(\alpha_i\) for every fixed \(i\).  If \(\alpha_i>0\), the vector
\(e_{n,i}\) exists for large \(n\) and
\[
 \mathfrak e_n(e_{n,i})=\nu_{n,i}
 =\frac{1}{M_n\alpha_{n,i}}\longrightarrow0.
\]
Assumption (i), and a further diagonal extraction over these indices,
give strong limits \(e_{n,i}\to e_i\).  If \(\alpha_i=0\), its rank-one
term has Hilbert--Schmidt norm \(\alpha_{n,i}\to0\) and needs no eigenvector
limit.  It follows that each fixed finite truncation
\[
 \sum_{i\le N}\alpha_{n,i}\,e_{n,i}\otimes e_{n,i}
\]
converges in Hilbert--Schmidt norm.  Here \(e\otimes e\) denotes the
rank-one operator \(F\mapsto\ip{F}{e}_H e\), and a term with zero
eigenvalue is omitted.  For the positive limiting coefficients this uses
\(\norm{e_{n,i}\otimes e_{n,i}-e_i\otimes e_i}_{\HS}
\le2\norm{e_{n,i}-e_i}_H\).

After division by \(M_n\ge1\), the squared tail is still at most \(r_N^2\)
by \eqref{eq:inverse-square-tail}.  The convergent finite truncations and
this uniform tail show that \(k_n/M_n\) has an \(L^2\) limit \(k^*\).
Hilbert--Schmidt convergence implies operator-norm convergence, so the
operator with kernel \(k^*\) has norm one.  On the other hand,
assumption (iii) implies \(\mathfrak D(k_n/M_n)\to0\) almost everywhere,
since its unnormalized numerator has a finite almost-everywhere limit.
Boundedness of \(\mathfrak D\) also gives convergence in \(L^2\) to
\(\mathfrak Dk^*\).  Both imply convergence in measure, whose limit is
unique.  Thus \(\mathfrak Dk^*=0\).  Centering passes to an \(L^2\) limit,
so \eqref{eq:double-difference-injective} gives \(k^*=0\), contradicting
its operator norm one.  This proves boundedness of \(M_n\).

Take now any subsequence without dividing by \(M_n\).  The bounded inverse
norms allow a numerical diagonal extraction
\(\kappa_{n,i}\to\kappa_i\) for every \(i\).  A positive \(\kappa_i\)
gives \(\mathfrak e_n(e_{n,i})=1/\kappa_{n,i}\) bounded, so assumption (i)
applies to its eigenvectors.  A further diagonal extraction over every
index with \(\kappa_i>0\) gives simultaneous strong limits
\(e_{n,i}\to e_i\).  A zero limiting coefficient again makes its rank-one
term vanish.  The same truncation
argument using \eqref{eq:inverse-square-tail} supplies a further
\(L^2\) limit \(k^*\) of the kernels.  It is doubly centered.  Assumption
(iii), boundedness of \(\mathfrak D\), and uniqueness of limits in measure
now give \(\mathfrak Dk^*=\mathfrak Dk\).  Equation
\eqref{eq:double-difference-injective} implies \(k^*=k\).  Every subsequence
has a further subsequence with this same \(L^2\) limit, proving the full
claimed convergence.
\end{proof}

We apply this criterion to the graph inverses.  Let \(K_n^V\) be the
inverse of \(B_n\) on centered vertex functions, extended by zero on
constants, and write \(k_n^V(x,y)\) for its kernel relative to \(\mu_n\).
Thus
\[
 (K_n^Vg)(x)=\sum_{y\in V_n} k_n^V(x,y)g(y)\mu_n(y).
\]
Its column equation follows from \(B_nK_n^V=\Id-\Pi_{\mathrm{const}}\)
and Lemma~\ref{lem:normalization}:
\[
 Q_n k_n^V(\,\cdot\,,y)=\delta_y-\mu_n,
 \qquad
 Q_n\bigl(k_n^V(\,\cdot\,,y)-k_n^V(\,\cdot\,,y')\bigr)
       =\delta_y-\delta_{y'}.
\]
Here \(\mu_n\) in the first equation is its vector of vertex masses.
In particular, a difference of columns is exactly a unit-current voltage
for the all-edge unit-conductance graph.  No factor involving a degree or
\(2n\) remains in this equation.

Define its kernel on the first tour by
\[
 k_n(t,s)=k_n^V\bigl(v_n^{(1)}(\lfloor2nt\rfloor),
                         v_n^{(1)}(\lfloor2ns\rfloor)\bigr)
\]
away from the irrelevant endpoints.  Its integral operator is precisely
the inverse on \(\mathcal H_n\cap\ind^\perp\) extended by zero on constants and on
\(\mathcal H_n^\perp\).  Indeed, integrating an arbitrary time function over the
occurrences of a vertex is its orthogonal projection to \(\mathcal H_n\), multiplied
by that vertex's mass \(\mu_n(v)\).  The kernel is symmetric and doubly
centered.

Let \(G\) be the round Green kernel for the unhalved Dirichlet integral in
Proposition~\ref{prop:continuum-form}, and let \(G_\mu\) be its centering in
both variables there.  This normalization agrees with the voltage in
\eqref{eq:round-voltage}: on the unit round sphere,
\[
 G(x,y)=-\frac1{2\pi}\log\operatorname{ch}(x,y)+c_0,
\]
where \(c_0\) gives zero round mean.  Indeed the round Laplacian of
\(\log\operatorname{ch}(x,y)\) is \(-1/2\) away from \(y\), and its
logarithmic singularity has flux \(2\pi\); hence the displayed function
has the distributional normalization \eqref{eq:round-green-normalization}.
Define
\[
 k(t,s)=G_\mu(\eta_1(t),\eta_1(s)).
\]
The proposition and the area pushforward show that \(k\) is doubly centered
and belongs to \(L^2([0,1]^2)\).  For almost every quadruple
\((t,t',s,s')\), the four points
\[
 x=\eta_1(t),\quad x'=\eta_1(t'),\quad y=\eta_1(s),\quad y'=\eta_1(s')
\]
are distinct: their joint image law is \(\mu^{\otimes4}\), and \(\mu\) is
nonatomic.  On such a quadruple, uniform tour tracking makes the four
corresponding vertex sequences converge to these points.
Theorem~\ref{thm:four-point} applies pathwise to these sequences.  Its
unit-current normalization and the column equation above yield
\begin{equation}\label{eq:voltage-kernel-limit}
 \mathfrak Dk_n(t,t',s,s')\longrightarrow
 G(x,y)-G(x,y')-G(x',y)+G(x',y')
 =\mathfrak Dk(t,t',s,s').
\end{equation}
Centering cancels from the displayed four-point combination.  This
argument only uses distinct limiting positions; it makes no estimate near
coincident positions.

Lemma~\ref{lem:energy-compactness} verifies assumptions (i) and (ii) of
Lemma~\ref{lem:inverse-kernel-compactness}, with \(p=7/10>1/2\), while
\eqref{eq:voltage-kernel-limit} verifies (iii).  We have proved
\begin{equation}\label{eq:inverse-HS-convergence}
 \norm{k_n-k}_{L^2([0,1]^2)}
 =\norm{K_n-K}_{\HS}\longrightarrow0,
\end{equation}
where \(K\) is the integral operator with kernel \(k\).  By
Lemma~\ref{lem:tour-order}, \(U_1\) is unitary; the pushforward formula
therefore identifies
\[
 K=U_1K_hU_1^{-1},
\]
where \(K_h\) is the centered inverse of the continuum operator \(B_h\)
from Proposition~\ref{prop:continuum-form}.  This also identifies the
constant mode.  The two parts of the argument have separate roles:
four-point voltages identify the only possible centered kernel limit, while
the block estimate supplies the compactness that ensures such a limit in
\(L^2\).

\subsection{Eigenvalues, heat traces, and the metric-measure coupling}

\begin{proof}[Proof of Theorem~\ref{thm:main}]
Start with any sequence of sizes tending to infinity and pass to a common
representation extraction as above.  On its probability-one event,
\eqref{eq:inverse-HS-convergence} holds.  Let \(\kappa_i\), \(i\ge1\), be
the positive eigenvalues of \(K\) in decreasing order, with multiplicity.
Proposition~\ref{prop:continuum-form} says that \(K_h\) is injective on
the centered subspace of \(L^2(\mu)\), which is infinite-dimensional.
Thus every \(\kappa_i\) is positive and
\[
 \kappa_i=\nu_i^{-1},\qquad \nu_i=2\Lambda_i(h),\quad i\ge1.
\]
For a nonnegative compact self-adjoint operator \(T\) on \(H\), let
\(\kappa_i(T)\) denote its positive eigenvalues in decreasing order, with
multiplicity and with zeros appended after a finite list.  The variational
formula is
\[
 \kappa_i(T)=
 \inf_{\substack{F\subset H\\ \dim F=i-1}}
 \ \sup_{\substack{u\in F^\perp\\ \norm{u}_H=1}}
       \ip{Tu}{u}_H,\qquad i\ge1,
\]
where \(F\) ranges over linear subspaces.  When \(\kappa_i(T)>0\), the
spectral theorem gives the upper bound by choosing the span of the first
\(i-1\) eigenvectors, and the lower bound because \(F^\perp\) meets the
span of the first \(i\) eigenvectors nontrivially for every such \(F\).
When \(\kappa_i(T)=0\), positivity gives the lower bound and a trial
subspace containing the positive range gives the upper bound.  For two
such operators \(T,S\), their quadratic forms on any unit vector differ
by at most \(\norm{T-S}_{\mathrm{op}}\).  Comparing the two suprema for
each \(F\), and then the two infima, therefore gives
\[
 \abs{\kappa_i(T)-\kappa_i(S)}\le\norm{T-S}_{\mathrm{op}}.
\]
Applying this to \(K_n,K\) yields
\[
 \abs{\kappa_{n,i}-\kappa_i}
 \le\norm{K_n-K}_{\mathrm{op}}
 \le\norm{K_n-K}_{\HS}\longrightarrow0
 \qquad(i\ge1).
\]
The characterization applies with multiplicities and with zeros appended
to the finite-rank spectra.  Since \(\kappa_i>0\), for each fixed \(i\)
the number \(\kappa_{n,i}\) is positive for all sufficiently large \(n\),
and reciprocation gives \(\nu_{n,i}\to\nu_i\).  The constant eigenvalue is
zero on both sides.  With the conventions of
Lemma~\ref{lem:normalization} and Proposition~\ref{prop:continuum-form},
\begin{equation}\label{eq:target-eigenvalue-convergence}
 n\lambda_{n,i}=\frac{\nu_{n,i}}2
 \longrightarrow\frac{\nu_i}2=\Lambda_i(h)
 \quad\text{for every }i\ge0.
\end{equation}
Zero padding for the inverses is exactly the prescribed infinite padding
for the finite generator spectra.  This proves coordinatewise convergence,
including every multiplicity, with acceleration exactly \(n\).

For the heat traces, use \(e^{-t\infty}=0\) to write
\[
 H_n(t)=1+\sum_{i\ge1}e^{-t\nu_{n,i}/2},\qquad
 H_h(t)=1+\sum_{i\ge1}e^{-t\nu_i/2}.
\]
Fix \(0<\delta<T<\infty\).  For \(N\ge i_0\), the high-index bound
\eqref{eq:raw-eigenvalue-tail} gives, for all sufficiently large \(n\),
\begin{equation}\label{eq:uniform-heat-tail}
 \sup_{\delta\le t\le T}\sum_{i>N}e^{-t\nu_{n,i}/2}
 \le\sum_{i>N}e^{-\delta c_* i^p/2}\longrightarrow0
 \quad\text{as }N\to\infty.
\end{equation}
The summability follows, for example, by comparison with the integral of
\(e^{-\delta c_*x^p/2}\).  Passing to the limit in each fixed index in
\eqref{eq:raw-eigenvalue-tail} gives the same bound for \(\nu_i\) and the
same estimate for the continuum tail.  For a fixed finite set of indices,
\eqref{eq:target-eigenvalue-convergence} gives uniform convergence of the
corresponding exponential terms on \([\delta,T]\), since for finite
\(x,y\ge0\),
\[
 \sup_{\delta\le t\le T}\abs{e^{-tx/2}-e^{-ty/2}}
 \le\frac T2\abs{x-y}.
\]
Combining the finite sums and \eqref{eq:uniform-heat-tail} proves
\(H_n\to H_h\) uniformly on \([\delta,T]\), hence in
\(C_{\mathrm{loc}}((0,\infty))\).  Finiteness and measurability of the
limiting heat trace were established in Proposition~\ref{prop:continuum-form}.

It remains to express the metric and measure convergence in the requested
intrinsic topology.  This uses the same coordinate sphere as the inverse
kernel above.  Write \(\Gamma_n\) and \(\Gamma\) for the two graphs in
\eqref{eq:common-distance-graph}, and let \(\varepsilon_n\to0\) bound their
Hausdorff distance in the maximum product metric of round distance and
ordinary distance on \([0,\infty)\).  By uniform continuity of \(D_h\) on
\(\Sph\times\Sph\), its modulus
\[
 \omega_h(r)=\sup\bigl\{\abs{D_h(x,y)-D_h(x',y')}:
 d_{\mathrm{round}}(x,x')+d_{\mathrm{round}}(y,y')\le r\bigr\}
\]
tends to zero as \(r\downarrow0\).  Approximating each point of
\(\Gamma_n\) by a point of \(\Gamma\) gives
\begin{equation}\label{eq:coordinate-distortion}
 \sup_{u,v\in V_n}
 \abs{a_nd_n(u,v)-D_h(z_n(u),z_n(v))}
 \le\varepsilon_n+\omega_h(2\varepsilon_n)\longrightarrow0.
\end{equation}
Approximating \((x,x,0)\in\Gamma\) in the opposite direction also shows
that \(z_n(V_n)\) is \(\varepsilon_n\)-dense in the round sphere.

The pushed measures \(\bar\mu_n=(z_n)_*\mu_n\) converge weakly to \(\mu\)
by Proposition~\ref{prop:geometry}.  On the compact round sphere this is
convergence in Prokhorov distance.  Choose \(\rho_n\downarrow0\) larger
than that distance and than \(\varepsilon_n\).  The coupling
characterization of Prokhorov distance gives a coupling of \(\bar\mu_n\)
and \(\mu\) assigning mass at least \(1-\rho_n\) to pairs at round
distance at most \(\rho_n\).  Lift it to a coupling of \(\mu_n\) and
\(\mu\) by the finite conditional distribution of a vertex given its
projection \(z_n(v)\).  The relation
\[
 \mathcal R_n=\{(v,x)\in V_n\times\Sph:
                d_{\mathrm{round}}(z_n(v),x)\le\rho_n\}
\]
is a correspondence, by round density, and the lifted coupling gives it
mass at least \(1-\rho_n\).  Its distortion for \(a_nd_n\) and \(D_h\)
is at most
\[
 \varepsilon_n+\omega_h(2\varepsilon_n)+\omega_h(2\rho_n),
\]
by \eqref{eq:coordinate-distortion}.  This tends to zero.  For completeness,
if a correspondence \(\mathcal R\) between spaces \(X,Y\) has distortion
less than \(2\xi\), the cross-distance
\[
 d(x,y)=\inf_{(x',y')\in\mathcal R}
          \bigl(d_X(x,x')+\xi+d_Y(y',y)\bigr)
\]
together with the original distances gives a metric on their disjoint
union.  The distortion inequality verifies the triangle inequalities that
cross between the spaces.  Related pairs have distance at most \(\xi\),
so the Hausdorff distance is at most \(\xi\).  A coupling giving the
relation mass at least \(1-\rho\) makes the Prokhorov distance at most
\(\max\{\xi,\rho\}\).  Apply this with
\(\xi=\tfrac12\operatorname{dis}(\mathcal R_n)+1/n\) and
\(\rho=\rho_n\).  We obtain
\[
 (V_n,a_nd_n,\mu_n)\longrightarrow(\Sph,D_h,\mu_h)
 \quad\text{in Gromov--Hausdorff--Prokhorov distance}.
\]
The normalization \(a_n\) here is exactly the deterministic diameter
normalization in Theorem~\ref{thm:main}; its identification with the scale
in the companion distance graph was proved in
Lemma~\ref{lem:space-normalization}.

All three convergences have been proved on this one representation with
the same \(\mu_h\), \(D_h\), and conformal sphere.  The conformal or
anticonformal sign retained by Proposition~\ref{prop:geometry} preserves
the round Dirichlet integral, so it also preserves the spectrum defined
from this sphere.  The continuum spectral readouts are measurable
functions of its area in that coordinate by
Proposition~\ref{prop:continuum-form}.  Thus the joint limiting law on
this extraction is the one in Theorem~\ref{thm:main}, with its stated clock.

Every original subsequence permits a further common extraction with these
almost-sure conclusions.  The Gromov--Hausdorff--Prokhorov space of compact
metric probability spaces, the coordinatewise product
\([0,\infty]^{\mathbb N_0}\), and \(C_{\mathrm{loc}}((0,\infty))\) are
Polish.  The subsequence criterion for convergence in distribution
therefore gives the asserted joint convergence through all positive
integers \(n\).
\end{proof}
